\SourceLeaf{357}

\noindent wenn man $X = m'm''$, $X' = mm''$, $X'' = mm'$ einsetzt, während der sechsfache
Rauminhalt der Pyramide, die sich dadurch mit dem Nullpunkt bildet, $= mm'm''\BareRoot{D}$
wird. Folglich [ist] das Lot
\[
= \BareRoot{\dfrac{D}{F\left(\dfrac{1}{m}, \dfrac{1}{m'}, \dfrac{1}{m''}\right)}};
\]

$T$, $T'$, $T''$ beziehen sich ebenso auf die Form $\begin{pmatrix} AD, A'D, A''D \\ BD, B'D, B''D \end{pmatrix}$
wie $t$, $t'$, $t''$ auf $\begin{pmatrix} a, a', a'' \\ b, b', b'' \end{pmatrix}$.



\begin{center}
\bfseries BEMERKUNG.
\end{center}

Die vorstehende Notiz befindet sich in einem Handbuch und stammt aus dem Juli 1831. Sie
zeigt, dass sich \textsc{Gauss} auch mit komplexen Größen von mehr als zwei Einheiten beschäftigte.
Solche Größen hatte er bereits in der Schlussbemerkung seiner Selbstanzeige der Theoria residuorum
biquadraticorum, Commentatio secunda vom 23. April 1831 (Bd. II, S. 169 bis 198) erwähnt.
Deshalb wurde die Notiz hier noch einmal aufgenommen, obwohl sie schon in Bd. II, S. 305 bis
307 wegen ihrer Beziehung zur Theorie der ternären Formen abgedruckt ist.
\begin{flushright}
\textsc{Stäckel.}
\end{flushright}


